% Seção: Responsabilidades
% Dono: Rafael · Orçamento: 0,5 página · docs/fontes/plano-de-projeto.md §14.1
% Conteúdo: quem fez o quê, a partir da tabela-mestre
% (docs/plano-de-desenvolvimento/README.md) e do git log.
% v1 (03/11): divisão planejada. No P3, confirmar com o git log e com as features que de
% fato fecharam o Definition of Done; o que saiu do escopo sai da tabela.
\section{Responsabilidades}\label{sec:responsabilidades}

A Tabela~\ref{tab:responsabilidades} resume a divisão do trabalho. Cada
integrante foi dono de uma frente: a fundação de uma área do código e as
funcionalidades que a estendem.

\begin{table}[ht]
\centering
\caption{Divisão do trabalho entre os integrantes.}
\label{tab:responsabilidades}
\footnotesize
\begin{tabularx}{\textwidth}{@{}l>{\raggedright\arraybackslash}X>{\raggedright\arraybackslash}X@{}}
\toprule
Integrante & Fundação e integração & Funcionalidades \\
\midrule
Thiago   & Núcleo de grafos, heap e buscas & Centralidades \\
Miguel   & Leitura do GTFS, modelo normalizado e rede de teste & Espera e simulação de trajeto \\
Henrique & Modelagem, diagramas UML e construtores das visões & Trajeto mais rápido \\
Arthur   & Aplicação Streamlit e componente de mapa & Eficiência da rede \\
Caio     & Repositório, verificação de imports e guia do projeto & Falhas e resiliência \\
Rafael   & Relatório e medição de tempo & Menos baldeações; avaliação de desempenho \\
\bottomrule
\end{tabularx}
\end{table}

\pendente{P3: acrescentar as desejáveis e opcionais que foram concluídas e conferir
com o \texttt{git log}.}
